%@ pages 36-49
%@ status translated
% Original section 3, "Les n-préchamps de Segal", pp. 36--49.

\clearpage
\section{Segal \texorpdfstring{$n$}{n}-prestacks}
\label{nchampsegpage}

Let $\Xx$ now be a site, i.e. a category equipped with a Grothendieck topology.
Recall that this means a collection of sieves $\Bb
\subset \Xx /X$, with $X$ varying, satisfying certain axioms
(see \cite{SGA4} \cite{MacLaneMoerdijk} \cite{Rezk}).


An {\em $n$-Segal prestack over $\Xx$} is a functor from
$\Xx ^o$ to the category $nSePC$ of Segal $n$-precategories, in other words
it is a presheaf of Segal $n$-precategories on $\Xx$. For example, for
$n=0$ it is simply a
presheaf of simplicial sets on $\Xx$.

We denote by $nSePCh(\Xx )$ or simply $nSePCh$ the category of
Segal $n$-prestacks over $\Xx$ with the morphisms strictly respecting the
structure (i.e. the morphisms of presheaves with values in
$nSePC$).

Let us insist on the fact that these are strict functors from $\Xx ^o$ to
the $1$-category $nSePC$ in which the morphisms are the morphisms of
Segal $n$-precategories strictly respecting the structure; in particular
we {\em
do not consider for the moment} any ``weak presheaves'' or
``lax'' ones.
This will be justified later by Theorem \ref{correlation}
and in
particular by its ``strictification'' part, which is in fact Corollary
\ref{strictif3}.

One may apply the operation $\tau _{\leq m}$ to a Segal $n$-prestack
$A$:
first
one replaces each $A(X)$ by the Segal $n$-category $SeCat(A(X))$ and
then one applies object by object the truncation $\tau _{\leq m}$ described
in the previous section.


For $n=0$ one has at one's disposal Jardine's closed model category \cite{Jardine}
of simplicial presheaves: the cofibrations are the
injections, the weak equivalences are Illusie's \cite{Illusie},
and the fibrations are determined by the lifting property.
Jardine shows in \cite{Jardine} that these three classes give a
structure of
closed model category.



For $n>0$ we shall
follow the thread of Jardine's argument \cite{Jardine} (which is in some sense
summarised in Lemma \ref{dhklemme}).
To begin with, we impose---as Jardine does---that the {\em cofibrations} of
Segal $n$-prestacks be the injections,
i.e. the morphisms $A\rightarrow B$ such that $A(X)\rightarrow B(X)$ is a
cofibration for every $X\in \Xx$.

The main difference with \cite{Jardine}
is that one must give a recursive definition (in $n>0$) of
weak equivalence. We shall say that a morphism $A\rightarrow B$ of Segal
$n$-prestacks,
where all
the $A(X)$ and $B(X)$ are Segal $n$-categories, is
a {\em weak equivalence} if:
\newline
---for every $X\in \Xx$ and all $x,y\in A_0(X)$ the morphism
$$
A_{1/}(x,y)\rightarrow B_{1/}(x,y)
$$
is a weak equivalence of Segal $n-1$-prestacks over
$\Xx /X$;
and
\newline
---the morphism of presheaves of sets
$\tau _{\leq 0} A\rightarrow \tau _{\leq 0}B$ on $\Xx$
induces a surjection on the associated sheaves.

We shall say that a morphism $A\rightarrow B$ of Segal $n$-prestacks is
a {\em
weak equivalence} if the morphism $SeCat(A)\rightarrow SeCat(B)$ is a
weak equivalence in the sense of the definition above.

Recall that for $n=0$ we have taken Illusie's equivalence as the notion
of weak
equivalence (the definitions above not really making sense).
One may observe that this choice is compatible with the construction $\Pi
_{m,Se}\circ A$ and with the truncation operation, see Lemma \ref{compatible}
and Corollary \ref{compPin} below.


A morphism of Segal $n$-prestacks is a {\em fibration} if it has
the lifting property with respect to
trivial cofibrations.


{\em Comparison with the coarse topology:}
Let us denote by $\Gg$ the Grothendieck topology given on $\Xx$ in order
to distinguish it from the
coarse topology.
Recall that the coarse topology is the one in which the only sieve
covering $X\in \Xx$ is
$\Bb = X/\Xx$ (in other words, for a family to cover $X$ it must
contain $X$).
We shall often use this topology, so it is worth comparing the main
notions for
$\Gg$ and for the coarse topology.  The class of cofibrations does
not depend on
$\Gg$. Let $f:A\rightarrow
B$ be a morphism of Segal $n$-prestacks. Then:
\newline
---$f$ is a
weak equivalence for the coarse topology if and only if for every
$X\in \Xx$, $A(X)\rightarrow B(X)$ is a weak equivalence of Segal $n$-precategories;
\newline
---if $f$ is a weak equivalence for the coarse topology then $f$
is a $\Gg$-weak equivalence;
\newline
---if $f$ is a $\Gg$-fibration then $f$ is a fibration for the coarse
topology;
\newline
---if $f$ is a fibration for the coarse topology then each $f_X:
A(X)\rightarrow B(X)$ is a fibration of Segal $n$-precategories;
\newline
---(here we use Theorem \ref{cmf} below) $f$ is a
trivial $\Gg$-fibration
if and only if $f$ is a trivial fibration for the coarse topology.

We now return to the consideration of our site $\Xx$ with its
topology $\Gg$.

\begin{theorem}
\label{cmf}
The category $nSePCh(\Xx )$ of Segal $n$-prestacks, with the three
classes of morphisms
(cofibrations, $\Gg$-equivalences and $\Gg$-fibrations), is a
closed model category.
\end{theorem}

For $n=0$ this is the theorem of
Jardine \cite{Jardine}, Joyal \cite{Joyal}, K.
Brown \cite{KBrown}. We
may therefore assume $n\geq 1$ and assume by induction
that the result is true for $n-1$.

For the proof of Theorem \ref{cmf} we shall need
the following notion. Let $f: A\rightarrow B$ be a morphism of
Segal $n$-prestacks.
We shall say that $f$ is a {\em preliminary equivalence} if:
\newline
---the morphism of presheaves of sets $A_0\rightarrow B_0$ induces an
isomorphism between the associated sheaves; and
\newline
---for every $X\in \Xx$ and for each sequence of objects $x_0,\ldots , x_m\in
A_0(X)$
the morphism
$$
A_{m/} (x_0,\ldots , x_m)\rightarrow B_{m/}(f(x_0),\ldots ,f(x_m))
$$
is a weak equivalence of Segal $n-1$-prestacks on $\Xx /X$.

\begin{lemma}
Let
$$
\begin{array}{ccc}
A & \stackrel{p}{\rightarrow} & B \\
\downarrow && \downarrow \\
A' & \stackrel{q}{\rightarrow} &B'
\end{array}
$$
be a commutative diagram of morphisms of Segal $n-1$-prestacks, where $B$ and
$B'$ are presheaves of sets. Suppose that the right-hand vertical morphism
(which we denote by $g$) induces an isomorphism between the associated
sheaves. Then
the morphism $A\rightarrow A'$ is a
weak equivalence of Segal $n-1$-prestacks if and only if
for every $X\in \Xx$ and $x\in B(X)$ the morphism
$p^{-1}(x)\rightarrow q^{-1}(g(x))$ is a weak equivalence of
Segal $n-1$-prestacks.
\end{lemma}
This is a direct consequence of the definition of weak equivalence.
\eop

\begin{corollary}
\label{etape}
A morphism $A\rightarrow B$ of Segal $n$-prestacks is a preliminary
equivalence if and only if for every $m\in \Delta$ the morphism
$A_{m/} \rightarrow B_{m/}$ of Segal $n-1$-prestacks is a weak
equivalence. In that case the morphisms
$$
A_{1/} \times _{A_0}\ldots \times _{A_0} A_{1/}\rightarrow
B_{1/} \times _{B_0}\ldots \times _{B_0} B_{1/}
$$
are also weak equivalences.
\end{corollary}
One notes that for $m=0$ this is the first condition for a
preliminary equivalence;  for the second condition one applies the
preceding lemma.
\eop

\begin{lemma}
\label{etape2}
Let $f: A\rightarrow B$ be a preliminary equivalence of Segal
$n$-prestacks.
Then $f$ is a weak equivalence of Segal $n$-prestacks.
\end{lemma}
{\em Proof:}
The case where the $A(X)$ and $B(X)$ are
Segal $n$-categories is immediate. It
therefore suffices to prove that if $f$ is a preliminary equivalence then
so is $SeCat(f): SeCat(A)\rightarrow SeCat(B)$. For this we shall
analyse the steps described in \cite{nCAT} for passing from $A$ to
$SeCat(A)$ (the discussion of \cite{nCAT} for the operation $Cat$ adapts
directly to the operation $SeCat$, as we recalled in \S 2).
Moreover, what is at issue here
is presheaves of Segal $n$-precategories on $\Xx$;
the operation $SeCat$ is performed object by object (it is functorial). The
construction $SeCat$ is obtained by iterating the constructions denoted
$Fix$ and
$Gen[m]$ in \cite{nCAT} and in \S 2 above.
The operation $Fix$ preserves the
type of weak equivalence of the $A_{m/}$, so it is clear that it
preserves the condition
of being a
preliminary equivalence. It therefore remains to deal with $Gen[m]$.

The operation $Gen[m]$ first comprises a step denoted $A\mapsto A'$ in
\cite{nCAT}, which preserves the type of weak equivalence of the $A_{p/}$. We
may therefore ignore this operation and again denote by $A$ (resp. $B$) its
result.

Next one considers a diagram of the form
$$
A_{m/} \rightarrow \Gg [m](A) \stackrel{g}{\rightarrow }
A_{1/}\times _{A_0} \ldots \times _{A_0} A_{1/}
$$
where $g$ is a trivial cofibration of Segal $n$-prestacks
object by object (i.e. for the coarse topology). (NB the notation
$\Gg [m]$ has no connection with the notation $\Gg$ for the Grothendieck
topology; we keep $\Gg [m]$ only in order to respect the
notations of \cite{nCAT}.) Now $Gen[m](A)_{p/}$ is the coproduct of
$A_{p/}$ and $A_{m/}\rightarrow \Gg [m](A)$ for various morphisms
$A_{m/}\rightarrow A_{p/}$ induced by various morphisms $p\rightarrow
m$.
The stability of $\Gg$-trivial cofibrations under
coproduct with cofibrant objects, for Segal $n-1$-prestacks,
is a
consequence of Theorem \ref{cmf} for $n-1$, via Reedy's lemma
\cite{Reedy}.
Hence the fact that the morphisms
$$
\Gg [m](A)\rightarrow \Gg [m](B)
$$
and
$$
A_{p/} \rightarrow B_{p/}
$$
are weak equivalences of Segal $n-1$-prestacks implies
that the same holds
for $Gen[m](A)_{p/}\rightarrow Gen[m](B)_{p/}$. We obtain from
Corollary \ref{etape} that $Gen[m](A)\rightarrow Gen[m](B)$ is a
preliminary equivalence. Observing that the weak equivalence of
Segal $n-1$-prestacks, and hence the preliminary equivalence of
Segal $n$-prestacks, is preserved by filtered colimits,
we obtain that
$SeCat(A)\rightarrow SeCat(B)$ is a preliminary equivalence,
hence a
weak equivalence.
\eop

\begin{lemma}
\label{locale}
The property of being a weak equivalence is local: if
$f:A\rightarrow
B$ is a morphism of $n$-prestacks on $\Xx /X$ and if $\Bb \subset
\Xx /X$
is a sieve covering $X$
then $f$ is a weak equivalence if and only if for every $Y\in
\Bb$, $f|_{\Xx /Y}$ is a weak equivalence.
\end{lemma}
{\em Proof:}
We argue by induction on $n$. For $n=0$ this is a direct consequence
of Illusie's definition. If we suppose that this is true for $n-1$
then it suffices (in view of the definition of weak equivalence) to observe
that the fact that $\tau _{\leq 0}SeCat(A)\rightarrow \tau
_{\leq 0} SeCat(B)$ induces
a surjection on the associated sheaves is a local property;
which is evident
(we use here the axioms of a Grothendieck topology).
\eop



{\em Proof of Theorem \ref{cmf}:}
Our proof is modelled on Jardine's argument \cite{Jardine}
and we therefore apply
our Lemma \ref{dhklemme}.

Property (0) is immediate.

It is easy to see that weak equivalences are stable under
retracts, by induction on $n$ (the case $n=0$ is a direct
consequence of Illusie's definition of weak equivalence).
The stability of cofibrations under retracts is immediate, which gives
condition (1) of
\ref{dhklemme}.



For the ``three for the price of two'' property ((2) of \ref{dhklemme}),
let $f: A\rightarrow B$ and $g: B\rightarrow C$ be two morphisms. It is
immediate that if $f$ and $g$ are
weak equivalences then so is $gf$; and that if $gf$ and $g$ are
weak equivalences then so is $f$ (one remarks that the first
condition for $g$ entails that the morphism $\tau _{\leq 0}B\rightarrow \tau
_{\leq 0}C$ induces an injection between the associated sheaves, so the
essential surjectivity of $gf$ implies that of $f$). It remains to show
that if $gf$
and $f$
are weak equivalences then $g$ is a weak equivalence.
The essential surjectivity condition is evident; the problem is the
first condition (see the analogous proof in \cite{nCAT}, Lemma
3.8). This condition is easy to verify for objects
$x,y\in B_0(X)$ of the form
$x=f(u)$ and $y=f(v)$ with $u,v\in A_0(X)$. For the general case, let
$x,y\in
B_0(X)$. Then for every $Y$ in a covering family of $X$ one can
find
objects $u_Y$ and $v_Y$ in $A_0(Y)$ with equivalences
$$
i: f(u_Y)\cong x|_Y, \;\;\; j: f(v_Y)\cong y|_Y
$$
(these are $1$-morphisms in $B(Y)$). One may say that ``composition with
$i$ and $j$'' induces an object-by-object equivalence between
$B_{1/}(f(u_Y), f(v_Y))|_{\Xx /Y}$ and
$B_{1/}(x,y)|_{\Xx /Y}$ (to make this precise one needs a discussion
analogous to that preceding Lemma 3.8 in \cite{nCAT}). The fact that
$$
B_{1/}(f(u_Y), f(v_Y))|_{\Xx /Y}\stackrel{\cong}{\rightarrow}
C_{1/}(gf(u_Y), gf(v_Y))|_{\Xx /Y}
$$
is an equivalence
implies that
$$
B_{1/}(x, y)|_{\Xx /Y}\rightarrow
C_{1/}(g(x), g(y))|_{\Xx /Y}
$$
is an equivalence. This for every $Y$ in the family covering $X$.
Since the condition of being a weak equivalence is local
(Lemma \ref{locale})
we obtain that
$$
B_{1/}(x, y)|_{\Xx /X}\rightarrow
C_{1/}(g(x), g(y))|_{\Xx /X}
$$
is an equivalence, which is the condition sought.

If $B'\rightarrow B$ is a morphism which has the lifting
property for every cofibration $A\rightarrow A'$, then one obtains in
particular the lifting property with respect to cofibrations of the form
$C_X\rightarrow C'_X$ where $C_X$ is the Segal $n$-prestack freely
generated by a Segal $n$-precategory $C$ over an object $X\in \Xx$;
and this morphism arises by the same construction from a cofibration
$C\rightarrow C'$. A morphism $C_X\rightarrow B'$ is identified with a
morphism $C\rightarrow B'(X)$, and hence for every $X\in \Xx$, $B'(X)\rightarrow
B(X)$ has the lifting property with respect to the cofibrations
of Segal $n$-precategories. This implies (by \ref{SeCmf}) that $B'(X)\rightarrow
B(X)$ is a weak equivalence, hence $B'\rightarrow B$ is a
weak equivalence object by object, i.e. for the coarse topology. One
thus obtains condition (3) of \ref{dhklemme}.

As usual we leave to the reader the set-theoretic questions
(4) and (5) of \ref{dhklemme}.


The stability of cofibrations under
coproduct ((6) of \ref{dhklemme}) is immediate.
It therefore only remains to prove the stability of trivial cofibrations
under coproduct, condition (7) of \ref{dhklemme}, which will occupy the rest of the
proof (we leave it to the reader to verify the stability of
trivial cofibrations under sequential limits).

Let $A\rightarrow B$ be a $\Gg$-trivial cofibration
and $A\rightarrow C$ a morphism of
Segal $n$-prestacks. Let $P:= B\cup ^AC$. We want to show that
$C\rightarrow P$ is a $\Gg$-trivial cofibration. It suffices to show
that $SeCat(C)\rightarrow SeCat(P)$ is a $\Gg$-weak equivalence; now we
know (Theorem \ref{SeCmf}) that the morphism
$$
SeCat(B)\cup ^{SeCat(A)}SeCat(C) \rightarrow SeCat(P)
$$
is a weak equivalence object by object. The morphism $SeCat(A)\rightarrow
SeCat(B)$ is still a $\Gg$-trivial cofibration. Up to replacing
$A$, $B$ and $C$ by $SeCat(A)$, $SeCat(B)$ and $SeCat(C)$ respectively, we may
therefore suppose that, for every $X$, $A(X)$, $B(X)$ and $C(X)$ are
Segal
$n$-categories.

We first treat a few particular cases.

\noindent
{\em Case 1.}\, We suppose that $A_0\rightarrow B_0$ induces an isomorphism
on the
associated sheaves. In this case it is easy to see that $A\rightarrow
B$ is
a preliminary equivalence. The result we are in the process of
proving,
applied to Segal $n-1$-prestacks and combined with Corollary
\ref{etape},
implies that $C\rightarrow B\cup ^AC$ is a preliminary equivalence.
Lemma \ref{etape2} then implies that it is a weak equivalence.


\noindent
{\em Case 2.}\, We suppose that $A\rightarrow B$ is a weak
equivalence and that for every
$X\in \Xx$ the morphism $A(X)\rightarrow B(X)$ is essentially
surjective.
We begin by remarking that the present Theorem \ref{cmf} is
immediate
in the case of the coarse
topology.
It follows that there exists a factorisation
$$
A\rightarrow B' \rightarrow B
$$
in which, for the coarse topology,
the first morphism is a cofibration and the second
is a trivial fibration.
Let $B'' \subset B'$ be defined by the condition that for every
$X\in \Xx$, $B''(X) \subset B'(X)$ is the full sub-Segal-$n$-category
whose objects are those in the image of $A_0$. The morphism
$B'' (X)\rightarrow B(X)$ is still essentially surjective and
fully faithful (since $B'(X)\rightarrow B(X)$ is an
equivalence of Segal $n$-categories). Hence $B'' \rightarrow B$ is
an equivalence for the coarse topology.
Our morphism factors
through a cofibration $A\rightarrow B''$.
Reedy's lemma \cite{Reedy}
(cf \cite{Hirschhorn} \cite{DHK} \cite{HoveyBook} \cite{JardineGoerssBook}),
which applies here since all objects are cofibrant, implies that the
morphism
$$
B'' \cup ^A C \rightarrow B\cup ^AC
$$
is an equivalence for the coarse topology. It then suffices to know
that
$$
C \rightarrow B'' \cup ^A C
$$
is a weak equivalence. The morphism $A_0\rightarrow B''_0$ is an
isomorphism (and $A\rightarrow B''$ is a weak equivalence since $B''$ is
object-by-object equivalent to $B$), so our Case 1 applies.

\medskip

\noindent
{\em End of the proof:}\,
We now treat the general case (with the same notation
$B\leftarrow A\rightarrow C$). Let $U_0\subset B_0$ be the image of $A_0$. Let
$U\subset B$ be the full sub-Segal-$n$-prestack having $U_0$ as
set of objects. Moreover,
for $X\in \Xx$, let $V_0(X)$ be the subset of the objects of $B_0(X)$
which are equivalent to an object of $U_0(X)$, and let $V\subset B$ be the
corresponding full sub-Segal-$n$-prestack. Object by object, the morphism
$U\rightarrow V$ is essentially surjective and fully
faithful, so our Case 2 applies.

The morphism $A\rightarrow U$ is a weak equivalence and it is bijective
on objects; hence Case 1 applies.

Finally the morphism of presheaves of sets $V_0\rightarrow B_0$ induces an
isomorphism between the associated sheaves: injectivity follows from the
fact that the
morphism of presheaves is injective; and surjectivity follows from the
fact that
every object of $B_0(X)$ is locally equivalent to an object which comes from
$A_0$, hence also from $V_0$. And so Case 1 applies once more.

We thus have a sequence of cofibrations
$$
A\rightarrow U \rightarrow V\rightarrow B
$$
along which the coproducts induce weak equivalences, by Cases 1, 2 and 1
respectively. This proves that the
coproduct along
$A\rightarrow B$ also induces a weak equivalence.
\eop



\subnumero{Compatibility with truncations}

Let $A$ be a Segal $m$-prestack, and $A\rightarrow A'$
an object-by-object equivalent replacement such that the $A'(X)$ are
Segal $m$-categories (e.g. $A'$ is a fibrant replacement for the
coarse
topology). One considers the $n$-prestack on $\Xx$:
$$
\tau _{\leq n} A := X\mapsto \tau _{\leq n}(A'(X)).
$$
We use here the word ``$n$-precategory'' in the sense of
\cite{nCAT} and ``$n$-prestack'' for presheaf of $n$-precategories.



{\em Remark:} this construction was denoted $\tau _{\leq n}^{\rm
pre}$ by the
second author in other papers; indeed, in those papers the notation
$\tau _{\leq n}$ had been reserved
for the stack associated (cf \S 9 and \S 13 below) to $\tau
_{\leq n}^{\rm
pre}$; but from the point of view adopted here, a prestack is weakly
equivalent to its associated stack, and one has no need to make this
distinction of notation.
One must however be careful that if $A$ is a stack (cf \S 9) then
$\tau _{\leq n}A$
is no longer necessarily a stack.



\begin{lemma}
\label{compatible}
Let $f:A\rightarrow B$ be a morphism of Segal $m$-prestacks. Then
$f$ is a weak equivalence for the topology $\Gg$ if and only if,
for every $n$,
$$
\tau _{\leq n} (A)\rightarrow
\tau _{\leq n} (B)
$$
is a $\Gg$-weak equivalence of $n$-prestacks.
\end{lemma}
{\em Proof:}
One first verifies that if $A$ and $B$ are $n$-prestacks of
$n$-groupoids
then $f:A\rightarrow B$ is a weak equivalence (using our
definition for Segal $n$-prestacks) if and only if $\Re
_{\geq 0} A
\rightarrow \Re _{\geq 0}B$ is an Illusie weak equivalence of
simplicial presheaves. We show this by induction on $n$, so we may
therefore suppose that it is true for $n-1$. For $x,y\in A_0(X)$ one has
$$
Path ^{x,y} \Re _{\geq 0} A \cong \Re _{\geq 0} (A_{1/}(x,y))
$$
(object by object one deduces this from \cite{Tamsamani}), and the same for
$B$. Thanks to the induction hypothesis we obtain that
$$
A_{1/}(x,y) \rightarrow B_{1/}(fx,fy)
$$
is a weak equivalence if and only if
$$
Path ^{x,y} \Re _{\geq 0} A \rightarrow Path ^{fx,fy} \Re _{\geq 0} B
$$
is an Illusie weak equivalence. On the other hand the presheaf
of sets
$\tau _{\leq 0} A$ is isomorphic to the presheaf of sets $\pi _0\circ
\Re _{\geq 0} A$, so we obtain that
$$
\tau _{\leq 0} A\rightarrow \tau _{\leq 0} B
$$
induces a surjection between the associated sheaves if and only if
$$
\pi _0\circ
\Re _{\geq 0} A\rightarrow \pi _0\circ
\Re _{\geq 0} B
$$
induces a surjection between the associated sheaves.
We conclude this part by remarking that a morphism $f:U\rightarrow V$ of
simplicial presheaves is an Illusie weak equivalence if and only
if (a) for every $x,y\in U_0(X)$, $Path ^{x,y}(U)\rightarrow Path^{fx,fy}(V)$
is an Illusie weak equivalence, and (b) the morphism $\pi _0\circ
U\rightarrow \pi _0\circ V$ induces a surjection between the associated sheaves
(the argument justifying this remark is immediate).

Now we treat the lemma for $m=0$: let $f: U\rightarrow V$ be a morphism
of simplicial presheaves. We know that the Postnikov truncation
$F \rightarrow\Re \tau _{\leq n}F$ respects
the $\pi _i$ for $i\leq n$. It follows that $f$ is an Illusie weak
equivalence if and only if
$$
\Re \tau _{\leq n}U\rightarrow \Re \tau _{\leq n}V
$$
is an Illusie weak equivalence for every $n$. Applying the
result of the preceding paragraph, we obtain that $f$ is
an Illusie weak equivalence
if and only if
$$
\tau _{\leq n}(f):\tau _{\leq n}U\rightarrow  \tau _{\leq
n}V
$$
is a weak equivalence for every $n$, which gives the lemma for $m=0$.

We now treat the general case by induction on $m\geq 1$; we
suppose
that the result is true for $m-1$.
For $f: A\rightarrow B$ and $x,y\in A(X)$ one has
$$
(\tau _{\leq n}A)_{1/}(x,y) \cong \tau _{\leq n-1}(A_{1/}(x,y))
$$
and likewise for $B$ (these equivalences entering into a commutative
diagram of naturality for $f$). The $A_{1/}(x,y)$ are
Segal $m-1$-prestacks. Using the induction
hypothesis, we obtain
that
$$
A_{1/}(x,y)\rightarrow B_{1/}(fx,fy)
$$
is an equivalence if and only if
$$
(\tau _{\leq n}A)_{1/}(x,y)\rightarrow (\tau _{\leq n}B)_{1/}(fx,fy)
$$
is an equivalence for every $n$.
On the other hand one has the natural isomorphism
$$
\tau _{\leq 0}A \cong \tau _{\leq 0}(\tau _{\leq n}A).
$$
Hence $\tau _{\leq 0}A\rightarrow \tau _{\leq 0}B$ induces a
surjection between the
associated sheaves if and only if
$$
\tau _{\leq 0}(\tau _{\leq n}A)\rightarrow \tau _{\leq 0}(\tau _{\leq n}B)
$$
induces a surjection between the associated sheaves for every $n$.
This completes the proof of the lemma.
\eop

\begin{corollary}
\label{compPin}
Let $m$ and $n$ be fixed. A morphism $f: A\rightarrow B$ of
Segal $m$-prestacks is
a $\Gg$-equivalence if and only if
$$
\Pi _{n,Se}\circ f: \Pi _{n,Se}\circ A\rightarrow \Pi _{n,Se}\circ A
$$
is a $\Gg$-equivalence of Segal $n+m$-prestacks.
\end{corollary}
\eop

\subnumero{Compatibility with the homotopy fibre product}

The closed model category above is almost proper. The only
problem is that $nSePC$ itself is not proper (for a counterexample,
see \cite{nCAT}). The $\Gg$-equivalences
are preserved
by fibre product with a fibration provided that the values
of the Segal $n$-prestacks considered are Segal
$n$-categories.
For $n=0$ (in which case there is no condition) the fact that the cmc of
simplicial presheaves is proper is proved by Jardine in
\cite{JardineBool}.

In fact one has a slightly stronger statement, namely the stability of
$\Gg$-equivalences under fibre product with a fibration for the coarse
topology. In other words, the homotopy fibre product for the coarse
topology preserves $\Gg$-equivalences. This statement was
conjectured for $n=0$ by C. Rezk in
an electronic letter to P. Hirschhorn \cite{Rezk}. Rezk
shows the analogous property
for presheaves of sets (i.e. the case of
$0$-stacks not of Segal). Even more interestingly, Rezk shows, in the framework
of localisation by an arbitrary pre-topology (i.e. collection of sieves),
that this property for ``local equivalences'' is
equivalent to the fact that the pre-topology is a topology, i.e.
satisfies the
axioms for a Grothendieck topology. The following lemma was therefore
inspired by this letter of Rezk and we thank P. Hirschhorn for
forwarding it to us.

\begin{lemma}
\label{proper}
If $A\rightarrow B$ is a cofibration and $A\rightarrow C$ a
$\Gg$-weak equivalence then $B\rightarrow B\cup ^AC$ is a
$\Gg$-weak
equivalence. On the other hand, if $f:F\rightarrow E$ is a fibration (even only
for the coarse topology) and $g:G\rightarrow E$ is a $\Gg$-weak
equivalence, and if $G(X)$ and $E(X)$ are Segal $n$-categories for every
$X\in
\Xx$ then $F\times _EG\rightarrow F$ is a $\Gg$-weak equivalence.
\end{lemma}
{\em Proof:}
The first part is immediate (it is the ``Reedy lemma''
\cite{Reedy} \cite{DHK} \cite{Hirschhorn}) since all objects are cofibrant.

For the second part, in general, we shall denote by
$F\times ^h_EG$ the object-by-object homotopy fibre product; it
is identified (object by object)
with the ordinary fibre product in the case where
$F\rightarrow E$ is a fibration for the
coarse topology and the $G(X)$ and $E(X)$ are Segal $n$-categories (see
\cite{nCAT} Theorem 6.7).
The truncation operation preserves---almost---the
homotopy fibre product, in the sense that one has an object-by-object
equivalence
$$
\tau _{\leq n-1} (F\times ^h _EG) \cong
\tau _{\leq n-1}\left( \tau _{\leq n}F\times ^h _{\tau _{\leq n}E}\tau
_{\leq n}G
\right) .
$$
This formula is proved for Segal $0$-prestacks
with the help of the exact sequence
of homotopy groups for the homotopy fibre product of spaces. One
then obtains it for Segal $n$-prestacks by induction on
$n$ (it is the
same argument as below).


Using this formula and Lemma
\ref{compatible}, one reduces to the case where $E$, $F$
and $G$ are $n$-truncated Segal $n$-prestacks whose values
are
Segal $n$-categories (they in fact correspond to $n$-prestacks not of
Segal). And we shall treat this case by induction on $n$. For example, if at the
outset one was in the case of simplicial presheaves, one reduces to the
case of simplicial presheaves which are object-by-object $n$-truncated.


We therefore suppose that $E,F,G$ are $n$-truncated Segal $n$-prestacks,
and that the values $E(X)$, $F(X)$ and $G(X)$ are Segal $n$-categories;
with a morphism $f: F\rightarrow E$ fibrant for the coarse topology,
and a morphism $g: G\rightarrow E$ which is a $\Gg$-weak equivalence.
The object-by-object homotopy fibre product is computed as the usual fibre
product $F\times _EG$. (Here we use Theorem 6.7 of
\cite{nCAT}
in its version adapted to Segal $n$-categories.)

If $(u,v)$ and $(u',v')$ are in $(F\times _EG)_0(X)$ then one has
(denoting by $w$ and $w'$ their images in $E_0(X)$):
$$
(F\times _EG)_{1/}((u,v), (u',v')) =
F_{1/}(u,u')\times _{E_{1/} (w,w')}G_{1/}(v,v').
$$
In this formula, on the right, the three terms are
$n-1$-truncated Segal $n-1$-prestacks on $\Xx /X$.
The structural morphisms are, on the one hand a fibration for the coarse
topology, and on the other hand a $\Gg$-weak equivalence.

 By induction
on $n$ we may suppose that the first projection
$$
(F\times _EG)_{1/}((u,v), (u',v'))
\rightarrow F_{1/}(u,u')
$$
is a $\Gg$-weak equivalence. On the other hand
$$
\tau _{\leq 0} (F\times ^h_EG) \rightarrow
(\tau _{\leq 0} F)\times _{\tau _{\leq 0} E} (\tau _{\leq 0} G)
$$
is a surjection of presheaves of sets.
The property in question for
presheaves of sets (see \cite{Rezk}) implies that the morphism
$$
(\tau _{\leq 0} F)\times _{\tau _{\leq 0} E} (\tau _{\leq 0} G)\rightarrow
\tau _{\leq 0} F
$$
induces a surjection on the associated sheaves; hence
$$
\tau _{\leq 0} (F\times ^h_EG)\rightarrow \tau _{\leq 0} F
$$
also induces
a surjection on the associated sheaves. We have verified everything
needed
to assert that $F\times ^h_EG\rightarrow F$ is a $\Gg$-weak
equivalence. \eop

\begin{corollary}
\label{cordereedy}
Let
$$
\begin{array}{ccccc}
F&\rightarrow & E & \leftarrow & G\\
\downarrow &&\downarrow && \downarrow \\
F'&\rightarrow & E' & \leftarrow & G'
\end{array}
$$
be a diagram in which the vertical morphisms are $\Gg$-weak
equivalences.
Then the morphism induced on the object-by-object homotopy product
$$
F\times ^h_EG\rightarrow F'\times ^h_{E'}G'
$$
is a $\Gg$-weak equivalence.
\end{corollary}
{\em  Proof:}
This is proved by a standard Reedy argument.
\eop

In \S 9 (Corollary \ref{rezkQ}) we shall see that the preceding corollary
implies that the
object-by-object homotopy fibre product is stable under the operation of
passage to the ``associated stack''. Rezk's question was in fact posed in this
form.

\subnumero{Descent data}



Let $Ho(nSePCh(\Xx ^{\rm gro}))$ denote the homotopy category
of the closed model category $nSePCh$ for $\Xx^{\rm gro}$.
Recall (\cite{Quillen}) that it may be defined as the Gabriel--Zisman
localisation \cite{GabrielZisman} of $nSePCh$ with respect to the
weak equivalences (for the coarse topology) or also as the
category whose objects are the Segal $n$-prestacks fibrant for the
coarse topology (and automatically cofibrant) and whose morphisms
are the equivalence classes of morphisms. In fact one has at one's disposal
the interval
$\overline{I}$ (the category with two isomorphic objects $0,1$) and, since all
objects are cofibrant, two morphisms $f,g:A\rightarrow B$ with $B$
fibrant are homotopic if and only if there exists $h: A\times
\overline{I}\rightarrow B$ with $h|_{A\times \{ 0\} }=f$ and
$h|_{A\times \{ 1\} }=g$.

Let $A, B\in nSePCh$. We denote by
$$
[A,B]_{\Xx ^{\rm gro}}
$$
the set of morphisms from $A$ to $B$ in $Ho(nSePCh(\Xx ^{\rm gro}))$.

In the same way one defines the homotopy category
 $Ho(nSePCh(\Xx ^{\Gg}))$ for the
topology $\Gg$,  either as the Gabriel--Zisman localisation of $nSePCh$ by
the $\Gg$-weak equivalences, or as the category of $\Gg$-fibrant objects
and homotopy classes of morphisms. We denote by
$$
[A,B]_{\Xx ^{\Gg}}
$$
the set of morphisms from $A$ to $B$ in $Ho(nSePCh(\Xx ^{\Gg}))$.

Let $X\in \Xx$ and $\Bb \subset \Xx /X$ a sieve covering $X$.
Let $A$ be a Segal $n$-prestack on $\Xx$. Let $\ast _{\Bb}$ be the
$n$-prestack on $\Bb$ with values the final Segal $n$-precategory
$\ast$. A {\em
descent datum for $A$ with respect to the sieve $\Bb$} is an element of
the set
$$
[\ast _{\Bb},A|_{\Bb}]_{\Bb ^{\rm gro}}.
$$
One may compute this set in the following way. Let $A\rightarrow A'$
be a fibrant replacement for the coarse topology. We shall see in the
next section that $A|_{\Bb}\rightarrow A'|_{\Bb}$ is a fibrant replacement
for $\Bb ^{\rm gro}$. Hence a descent datum, i.e.
an element of
$[\ast _{\Bb},A|_{\Bb}]_{\Bb ^{\rm gro}}$, is a morphism $\ast
_{\Bb}\rightarrow A'|_{\Bb}$, up to equivalence (the equivalence being
the homotopy relation with the interval $\overline{I}$, for example).
By a slight abuse of notation we shall say that a descent datum is a
morphism $\ast _{\Bb}\rightarrow A'|_{\Bb}$.

With the functor $p: \Bb \rightarrow \Xx$ and the construction $p_!$ (see \S 4
below) one may identify a morphism
$\ast _{\Bb}\rightarrow A'|_{\Bb}$ with a morphism
$p_!\ast _{\Bb}\rightarrow A'$. The Segal $n$-precategory $p_!\ast
_{\Bb}$
on $\Xx$ (which is in fact a presheaf of sets on $\Xx$) will
also often
be denoted
$\ast _{\Bb}$, and with this notation a descent datum for $A$ with
respect to $\Bb$ is a morphism $\ast _{\Bb} \rightarrow A'$ up to
homotopy equivalence.




For the question of ``the effectivity of descent data'' one is
interested
in descent data up to homotopy; that is why one uses
the homotopy category here to give the definition of
``the set of descent data''. One could also
envisage defining the {\em Segal $n$-category of descent
data}
of which the present set is the
$\tau _{\leq 0}$
%
%
\footnote{
The Segal $n$-category of descent
data would be $\underline{Hom}(\ast _{\Bb}, A')$ with the
notations of \S 11, which is equivalent
to  $\lim _{\leftarrow , \Bb}
A|_{\Bb}$ by \ref{calclim}.
The set of descent data
may then be expressed as
$$
[\ast _{\Bb},A|_{\Bb}]_{\Bb ^{\rm gro}}=\tau _{\leq 0} \lim _{\leftarrow , \Bb}
A|_{\Bb}.
$$};
%
%
in which case the term ``descent datum''
introduced here should be replaced by ``descent datum up to
homotopy''. However, for the use that will be made of it
in the present paper, the terminology given here is well suited.

One now has the equality
$$
[\ast _{\Xx /X},A|_{\Xx /X}]_{(\Xx /X) ^{\rm gro}}
= [\ast , A(X)],
$$
where the notation on the right means the set of homotopy classes of
morphisms $\ast
\rightarrow A(X)$ in $nSePC$; and the restriction of $\Xx /X$ to $\Bb$
therefore induces a morphism
$$
[\ast , A(X)]\rightarrow
[\ast _{\Bb},A|_{\Bb}]_{\Bb ^{\rm gro}}.
$$
We shall say that a descent datum for $A$ with respect to $\Bb$
is {\em effective} if it lies in the image of this morphism, i.e. if it
comes from an element of $[\ast , A(X)]$.
The notion of {\em effectivity of descent data} will enter into our
criterion for being a stack, \ref{critere} below.

In principle, descent data may arrive in the form of a diagram
$$
\ast _{\Bb} \leftarrow D \rightarrow A'\leftarrow A
$$
where the first and the last arrow are object-by-object weak equivalences.
In fact, descent data will most often come from diagrams of the form
$$
\ast _{\Bb} \leftarrow D \rightarrow A.
$$
In \S 5 below we shall construct a single arrow $D\rightarrow \ast
_{\Bb}$ independent of $A$, such that for every $n$-prestack $A$ whose
values $A(X)$ are fibrant Segal $n$-categories, the descent data for
$A$ come uniformly from diagrams
$\ast _{\Bb} \leftarrow D \rightarrow
A$.

One could also consider {\em generalised descent
data}, which one would define as elements of
$[\ast _{\Xx
/X}, A|_{\Xx /X}]_{\Xx ^{\Gg}}$.
In this optic such data may come, for example, from a diagram
$$
\ast _{\Xx /X} \leftarrow D \rightarrow A
$$
where the first arrow is a $\Gg$-weak equivalence. Such is
certainly the case if $D$ is object-by-object equivalent to $\ast_{\Bb}$
for a sieve $\Bb \subset \Xx /X$ covering $X$; but such a diagram
might also arrive from a {\em hypercovering} of $X$
(cf Verdier's exposé in \cite{SGA4b}).
We leave it to the reader to spell out the construction of an adequate $D$
from a hypercovering.
